% =====================================================================
%  Chapitre 10 : Authentification
% =====================================================================
\chapter{L'authentification des utilisateurs}
\label{ch:authentification}

\chapitreresume{Un catalogue public se consulte sans compte. Mais emprunter un
livre, réserver un ouvrage ou accéder à son historique demande une identité. Dans
ce chapitre, vous utilisez le système d'authentification de Django : inscription,
connexion, déconnexion, profil, protection des pages et gestion des sessions.}

\begin{objectifs}
\begin{itemize}
  \item Comprendre le modèle \code{User} fourni par Django.
  \item Activer les URL d'authentification intégrées.
  \item Créer les gabarits de connexion et de déconnexion.
  \item Écrire un formulaire d'inscription relié au profil \code{Membre}.
  \item Utiliser \code{login\_required} et \code{LoginRequiredMixin}.
  \item Rediriger l'utilisateur vers la page demandée après connexion.
  \item Configurer les sessions et les mots de passe.
  \item Afficher le profil et l'historique des emprunts.
\end{itemize}
\end{objectifs}

\section{Le système d'authentification de Django}

Django fournit l'application \code{django.contrib.auth}, déjà présente dans
\code{INSTALLED\_APPS}. Elle gère :
\begin{itemize}
  \item les utilisateurs, avec \code{User} ;
  \item les mots de passe, hachés avec PBKDF2 par défaut ;
  \item les sessions, stockées en base ;
  \item les groupes et les permissions ;
  \item les vues de connexion, de déconnexion et de réinitialisation.
\end{itemize}

\begin{retenir}
Django ne stocke jamais un mot de passe en clair. Il enregistre une empreinte
calculée avec un sel. La vérification recalcule l'empreinte et compare. Même un
administrateur de la base ne peut pas lire le mot de passe d'origine.
\end{retenir}

\subsection{Le modèle User}

\begin{table}[htbp]
\centering
\small
\begin{tabularx}{\textwidth}{L{4.4cm} X}
\toprule
\textbf{Champ} & \textbf{Rôle} \\
\midrule
\code{username} & Identifiant unique. Obligatoire par défaut. \\
\code{password} & Empreinte du mot de passe, jamais le mot de passe. \\
\code{email} & Adresse électronique, non unique par défaut. \\
\code{first\_name} & Prénom. \\
\code{last\_name} & Nom. \\
\code{is\_active} & Compte autorisé à se connecter. \\
\code{is\_staff} & Accès à l'administration. \\
\code{is\_superuser} & Tous les droits. \\
\code{date\_joined} & Date de création du compte. \\
\code{last\_login} & Dernière connexion réussie. \\
\bottomrule
\end{tabularx}
\caption{Les champs principaux du modèle User.}
\end{table}

\begin{note}
Pour un nouveau projet, Django recommande un modèle utilisateur personnalisé défini
avant la première migration. Dans ce module, nous gardons le modèle standard : il
suffit pour apprendre et pour le projet \gesibiblio. Le profil \code{Membre} ajoute
les informations propres à la bibliothèque.
\end{note}

\section{Activer les URL d'authentification}

Django fournit un jeu de vues prêt à l'emploi. Ajoutez cette ligne dans
\code{gesibiblio/urls.py} :

\begin{codebox}[python]{gesibiblio/urls.py}
from django.contrib import admin
from django.contrib.auth import views as auth_views
from django.urls import include, path

urlpatterns = [
    path("admin/", admin.site.urls),
    path("comptes/", include("django.contrib.auth.urls")),
    path("", include("catalogue.urls")),
]
\end{codebox}

\begin{note}
\code{django.contrib.auth.urls} n'est pas un module de vues : c'est une liste
d'URL. L'import \code{auth\_views} n'est utile que si vous personnalisez une vue.
Vous pouvez retirer la ligne d'import dans la version minimale.
\end{note}

Les routes créées sont :

\begin{table}[htbp]
\centering
\small
\begin{tabularx}{\textwidth}{L{5.4cm} L{5.0cm} X}
\toprule
\textbf{URL} & \textbf{Nom} & \textbf{Rôle} \\
\midrule
\url{/comptes/login/} & \code{login} & Connexion. \\
\url{/comptes/logout/} & \code{logout} & Déconnexion. \\
\url{/comptes/password_change/} & \code{password\_change} & Changement de mot de
passe. \\
\url{/comptes/password_reset/} & \code{password\_reset} & Demande de
réinitialisation. \\
\url{/comptes/password_reset/done/} & \code{password\_reset\_done} & Confirmation
d'envoi. \\
\url{/comptes/password_reset/<uidb64>/<token>/} &
\code{password\_reset\_confirm} & Saisie du nouveau mot de passe. \\
\url{/comptes/password_reset/complete/} & \code{password\_reset\_complete} &
Succès. \\
\bottomrule
\end{tabularx}
\caption{URL d'authentification fournies par Django.}
\end{table}

\begin{attention}
Sans les gabarits correspondants, ces URL provoquent une erreur
\code{TemplateDoesNotExist}. Django cherche les fichiers dans
\code{registration/} : \code{registration/login.html},
\code{registration/logged\_out.html}, etc.
\end{attention}

\section{Les gabarits de connexion}

\begin{codebox}[html]{templates/registration/login.html}
{% extends "base.html" %}

{% block titre %}Connexion | GesiBiblio{% endblock %}

{% block contenu %}
  <section class="formulaire">
    <h1>Connexion</h1>
    <p>Accédez à votre espace membre.</p>

    {% if form.errors %}
      <p class="erreur">
        Identifiant ou mot de passe incorrect. Réessayez.
      </p>
    {% endif %}

    <form method="post">
      {% csrf_token %}

      <div class="champ">
        <label for="{{ form.username.id_for_label }}">
          Nom d'utilisateur
        </label>
        {{ form.username }}
      </div>

      <div class="champ">
        <label for="{{ form.password.id_for_label }}">Mot de passe</label>
        {{ form.password }}
      </div>

      <input type="hidden" name="next" value="{{ next }}">

      <button class="bouton bouton-principal" type="submit">
        Se connecter
      </button>
    </form>

    <p>
      Pas encore de compte ?
      <a href="{% url 'catalogue:inscription' %}">Créer un compte</a>
    </p>
  </section>
{% endblock %}
\end{codebox}

\begin{attention}
Le message d'erreur reste volontairement vague : \og identifiant ou mot de passe
incorrect \fg. Dire \og cet utilisateur n'existe pas \fg{} permettrait à un
attaquant de découvrir les comptes valides. C'est une bonne pratique de sécurité.
\end{attention}

\subsection{La déconnexion}

Depuis Django 5.0, \code{LogoutView} n'accepte que \code{POST}. Un simple lien
renvoie une erreur \code{405}. Il faut un petit formulaire.

\begin{codebox}[html]{templates/partials/navbar.html}
<header class="entete">
  <nav class="conteneur barre">
    <a class="marque" href="{% url 'catalogue:accueil' %}">GesiBiblio</a>

    <button class="burger" type="button" aria-expanded="false"
            aria-controls="menu-principal">Menu</button>

    <ul class="menu" id="menu-principal">
      <li><a href="{% url 'catalogue:livre_liste' %}">Catalogue</a></li>

      {% if user.is_authenticated %}
        <li>
          <a href="{% url 'catalogue:profil' %}">
            {{ user.get_full_name|default:user.username }}
          </a>
        </li>
        <li>
          <form method="post" action="{% url 'logout' %}">
            {% csrf_token %}
            <button class="lien" type="submit">Déconnexion</button>
          </form>
        </li>
      {% else %}
        <li><a href="{% url 'login' %}">Connexion</a></li>
        <li>
          <a class="bouton bouton-principal"
             href="{% url 'catalogue:inscription' %}">Inscription</a>
        </li>
      {% endif %}
    </ul>
  </nav>
</header>
\end{codebox}

\begin{codebox}[html]{templates/registration/logged\_out.html}
{% extends "base.html" %}

{% block titre %}Déconnecté | GesiBiblio{% endblock %}

{% block contenu %}
  <h1>Vous êtes déconnecté.</h1>
  <p>À bientôt dans la bibliothèque.</p>
  <a class="bouton bouton-principal"
     href="{% url 'catalogue:accueil' %}">Retour à l'accueil</a>
{% endblock %}
\end{codebox}

\section{Les réglages d'authentification}

Ajoutez ces lignes dans \code{settings.py} :

\begin{codebox}[python]{gesibiblio/settings.py, extrait}
LOGIN_URL = "login"
LOGIN_REDIRECT_URL = "catalogue:profil"
LOGOUT_REDIRECT_URL = "catalogue:accueil"
\end{codebox}

\begin{description}[leftmargin=4.6cm,style=nextline]
  \item[\code{LOGIN\_URL}] Nom de la route de connexion. Utilisé par
  \code{login\_required} pour construire \code{/comptes/login/?next=...}.
  \item[\code{LOGIN\_REDIRECT\_URL}] Page d'arrivée après connexion, si aucun
  paramètre \code{next} n'est fourni.
  \item[\code{LOGOUT\_REDIRECT\_URL}] Page d'arrivée après déconnexion.
\end{description}

\begin{note}
Les valeurs sont des noms de route, pas des chemins. Django accepte aussi un chemin
littéral, mais le nom résiste aux changements d'URL. Pour une route d'application,
utilisez le nom complet \code{catalogue:profil}.
\end{note}

\section{Inscription et profil Membre}

Le modèle \code{Membre} est lié à \code{User} par une relation un à un. À
l'inscription, il faut créer les deux objets.

\subsection{Le formulaire d'inscription}

\begin{codebox}[python]{catalogue/forms.py, extrait}
from django import forms
from django.contrib.auth.forms import UserCreationForm
from django.contrib.auth.models import User


class InscriptionForm(UserCreationForm):
    email = forms.EmailField(
        required=True,
        label="Adresse électronique",
    )
    first_name = forms.CharField(
        max_length=80,
        required=True,
        label="Prénom",
    )
    last_name = forms.CharField(
        max_length=80,
        required=True,
        label="Nom",
    )
    telephone = forms.CharField(
        max_length=20,
        required=False,
        label="Téléphone",
    )

    class Meta(UserCreationForm.Meta):
        model = User
        fields = (
            "username",
            "first_name",
            "last_name",
            "email",
        )

    def clean_email(self):
        email = self.cleaned_data["email"].lower()
        if User.objects.filter(email__iexact=email).exists():
            raise forms.ValidationError(
                "Un compte utilise déjà cette adresse électronique."
            )
        return email
\end{codebox}

\begin{description}[leftmargin=5.1cm,style=nextline]
  \item[\code{UserCreationForm}] Fournit les champs \code{username},
  \code{password1} et \code{password2}, ainsi que la vérification de
  correspondance.
  \item[\code{clean\_email()}] Ajoute une unicité que le modèle \code{User}
  n'impose pas.
  \item[\code{lower()}] Normalise l'adresse pour éviter les doublons de casse.
\end{description}

\subsection{La vue d'inscription}

\begin{codebox}[python]{catalogue/views.py, extrait}
from django.contrib import messages
from django.contrib.auth import login
from django.shortcuts import redirect, render

from .forms import InscriptionForm
from .models import Membre


def inscription(request):
    if request.method == "POST":
        form = InscriptionForm(request.POST)
        if form.is_valid():
            user = form.save()
            Membre.objects.create(
                user=user,
                numero=f"M{user.pk:05d}",
                telephone=form.cleaned_data.get("telephone", ""),
            )
            login(request, user)
            messages.success(
                request,
                "Bienvenue ! Votre compte membre est créé.",
            )
            return redirect("catalogue:profil")
    else:
        form = InscriptionForm()

    return render(
        request,
        "catalogue/inscription.html",
        {"form": form},
    )
\end{codebox}

\begin{description}[leftmargin=4.9cm,style=nextline]
  \item[\code{form.save()}] Crée l'utilisateur et hache le mot de passe.
  \item[\code{Membre.objects.create}] Crée le profil lié, avec un numéro
  \code{M00042}.
  \item[\code{login(request, user)}] Ouvre la session immédiatement après
  l'inscription.
  \item[\code{f"M{user.pk:05d}"}] Formate l'identifiant sur cinq chiffres.
\end{description}

\begin{attention}
La création de l'utilisateur et celle du profil devraient être atomiques. Si la
création du \code{Membre} échoue, l'utilisateur reste sans profil. Entourez les deux
opérations d'un bloc \code{transaction.atomic()} :
\end{attention}

\begin{codebox}[python]{Création atomique}
from django.db import transaction

with transaction.atomic():
    user = form.save()
    Membre.objects.create(
        user=user,
        numero=f"M{user.pk:05d}",
    )
\end{codebox}

\begin{codebox}[html]{templates/catalogue/inscription.html}
{% extends "base.html" %}

{% block titre %}Inscription | GesiBiblio{% endblock %}

{% block contenu %}
  <h1>Créer un compte</h1>

  <form method="post" class="formulaire" novalidate>
    {% csrf_token %}

    {% if form.non_field_errors %}
      <div class="erreur">{{ form.non_field_errors }}</div>
    {% endif %}

    {% for champ in form %}
      <div class="champ">
        <label for="{{ champ.id_for_label }}">{{ champ.label }}</label>
        {{ champ }}
        {% if champ.help_text %}
          <p class="aide">{{ champ.help_text|safe }}</p>
        {% endif %}
        {% for erreur in champ.errors %}
          <p class="erreur">{{ erreur }}</p>
        {% endfor %}
      </div>
    {% endfor %}

    <button class="bouton bouton-principal" type="submit">
      Créer mon compte
    </button>
  </form>
{% endblock %}
\end{codebox}

\begin{note}
\code{help\_text|safe} est utilisé parce que \code{UserCreationForm} fournit des
aides contenant des balises \code{<ul>}. Ce contenu vient de Django, pas de
l'utilisateur. Ne faites jamais cela avec un texte saisi par un visiteur.
\end{note}

\subsection{La page de profil}

\begin{codebox}[python]{catalogue/views.py, extrait}
from django.contrib.auth.decorators import login_required
from django.shortcuts import get_object_or_404, render

from .models import Membre


@login_required
def profil(request):
    membre = get_object_or_404(Membre, user=request.user)
    emprunts = (
        membre.emprunts.select_related("livre")
        .order_by("-date_emprunt")
    )
    return render(
        request,
        "catalogue/profil.html",
        {"membre": membre, "emprunts": emprunts},
    )
\end{codebox}

\begin{codebox}[html]{templates/catalogue/profil.html}
{% extends "base.html" %}

{% block titre %}Mon profil | GesiBiblio{% endblock %}

{% block contenu %}
  <h1>{{ user.get_full_name|default:user.username }}</h1>

  <dl class="fiche-livre">
    <dt>Numéro de membre</dt>
    <dd>{{ membre.numero }}</dd>
    <dt>Adresse électronique</dt>
    <dd>{{ user.email }}</dd>
    <dt>Téléphone</dt>
    <dd>{{ membre.telephone|default:"Non renseigné" }}</dd>
    <dt>Inscrit le</dt>
    <dd>{{ membre.date_inscription|date:"d/m/Y" }}</dd>
  </dl>

  <h2>Mes emprunts</h2>

  <table class="tableau">
    <thead>
      <tr>
        <th>Livre</th>
        <th>Emprunt</th>
        <th>Retour prévu</th>
        <th>Statut</th>
      </tr>
    </thead>
    <tbody>
      {% for emprunt in emprunts %}
        <tr>
          <td>
            <a href="{{ emprunt.livre.get_absolute_url }}">
              {{ emprunt.livre.titre }}
            </a>
          </td>
          <td>{{ emprunt.date_emprunt|date:"d/m/Y" }}</td>
          <td>{{ emprunt.date_retour_prevue|date:"d/m/Y" }}</td>
          <td>
            {% if emprunt.date_retour %}
              Rendu le {{ emprunt.date_retour|date:"d/m/Y" }}
            {% elif emprunt.est_en_retard %}
              <span class="badge badge-retard">
                En retard de {{ emprunt.duree_retard }} jour{{
                emprunt.duree_retard|pluralize }}
              </span>
            {% else %}
              <span class="badge badge-disponible">En cours</span>
            {% endif %}
          </td>
        </tr>
      {% empty %}
        <tr>
          <td colspan="4">Aucun emprunt pour le moment.</td>
        </tr>
      {% endfor %}
    </tbody>
  </table>
{% endblock %}
\end{codebox}

\begin{attention}
\code{get\_object\_or\_404(Membre, user=request.user)} suppose que chaque
utilisateur a un profil. Un superutilisateur créé avec \code{createsuperuser} n'en
a pas. Soit vous créez le profil automatiquement, soit vous gérez l'absence avec
\code{hasattr(request.user, "membre")}.
\end{attention}

\section{Protéger une vue}

\subsection{Le décorateur login\_required}

\begin{codebox}[python]{Protéger une vue fonction}
from django.contrib.auth.decorators import login_required


@login_required
def profil(request):
    ...
\end{codebox}

Si l'utilisateur n'est pas connecté, Django le redirige vers
\code{LOGIN\_URL} avec \code{?next=/comptes/profil/}. Après connexion, la vue
\code{LoginView} lit \code{next} et renvoie l'utilisateur à la page demandée.

\begin{attention}
Ne faites pas confiance à \code{next} sans contrôle. Django fournit
\code{django.utils.http.url\_has\_allowed\_host\_and\_scheme} pour éviter une
redirection ouverte vers un site externe. Les vues intégrées appliquent déjà cette
vérification. Si vous écrivez votre propre redirection, faites de même.
\end{attention}

\subsection{Le mixin des vues de classe}

\begin{codebox}[python]{Protéger une vue classe}
from django.contrib.auth.mixins import LoginRequiredMixin
from django.views.generic import UpdateView

from .models import Membre


class ProfilUpdateView(LoginRequiredMixin, UpdateView):
    model = Membre
    fields = ["telephone", "adresse"]
    template_name = "catalogue/profil_modifier.html"

    def get_object(self, queryset=None):
        return self.request.user.membre
\end{codebox}

\begin{attention}
\code{LoginRequiredMixin} doit être placé à gauche dans l'héritage, avant la vue
générique. L'ordre des classes compte : les méthodes de la première classe
prennent le dessus. Un mixin placé à droite peut être ignoré.
\end{attention}

\section{Les sessions}

Une session associe un identifiant stocké dans un cookie à des données conservées
côté serveur. Django crée la session à la connexion.

\begin{codebox}[python]{Manipuler la session}
request.session["derniere_recherche"] = "prince"
request.session.get("derniere_recherche")
request.session.set_expiry(60 * 60 * 24 * 7)   # 7 jours
request.session.flush()                        # tout effacer
\end{codebox}

\begin{table}[htbp]
\centering
\small
\begin{tabularx}{\textwidth}{L{5.4cm} X}
\toprule
\textbf{Réglage} & \textbf{Effet} \\
\midrule
\code{SESSION\_COOKIE\_AGE} & Durée de vie en secondes. \code{1209600} par défaut,
soit deux semaines. \\
\code{SESSION\_EXPIRE\_AT\_BROWSER\_CLOSE} & Ferme la session à la fermeture du
navigateur. \\
\code{SESSION\_COOKIE\_SECURE} & Cookie uniquement en HTTPS. À activer en
production. \\
\code{SESSION\_COOKIE\_HTTPONLY} & Interdit l'accès au cookie en JavaScript.
Activé par défaut. \\
\code{SESSION\_SAVE\_EVERY\_REQUEST} & Écrit la session à chaque requête. Coûteux,
à laisser à \code{False}. \\
\bottomrule
\end{tabularx}
\caption{Les réglages de session les plus utiles.}
\end{table}

\begin{bonne}
En production, activez \code{SESSION\_COOKIE\_SECURE}, \code{CSRF\_COOKIE\_SECURE}
et \code{SECURE\_SSL\_REDIRECT}. Ces réglages empêchent l'envoi des cookies en
clair sur une connexion HTTP.
\end{bonne}

\section{Les validateurs de mot de passe}

Django applique quatre validateurs par défaut :

\begin{codebox}[python]{gesibiblio/settings.py, extrait}
AUTH_PASSWORD_VALIDATORS = [
    {
        "NAME": (
            "django.contrib.auth.password_validation"
            ".UserAttributeSimilarityValidator"
        ),
    },
    {
        "NAME": (
            "django.contrib.auth.password_validation"
            ".MinimumLengthValidator"
        ),
        "OPTIONS": {"min_length": 9},
    },
    {
        "NAME": (
            "django.contrib.auth.password_validation"
            ".CommonPasswordValidator"
        ),
    },
    {
        "NAME": (
            "django.contrib.auth.password_validation"
            ".NumericPasswordValidator"
        ),
    },
]
\end{codebox}

\begin{description}[leftmargin=5.4cm,style=nextline]
  \item[\code{UserAttributeSimilarityValidator}] Refuse un mot de passe trop proche
  du nom ou de l'adresse.
  \item[\code{MinimumLengthValidator}] Impose une longueur minimale. Neuf
  caractères est un plancher raisonnable.
  \item[\code{CommonPasswordValidator}] Refuse les mots de passe présents dans une
  liste de plusieurs milliers de valeurs courantes.
  \item[\code{NumericPasswordValidator}] Refuse un mot de passe uniquement
  numérique.
\end{description}

\begin{note}
Une longueur minimale de neuf caractères protège mieux qu'une exigence de symboles
compliqués, que les utilisateurs contournent en ajoutant \code{!} à la fin. Les
phrases de passe longues sont préférables.
\end{note}

\section{Réinitialiser un mot de passe}

Les vues sont déjà disponibles. Il faut fournir les gabarits et configurer l'envoi
de courrier.

\begin{codebox}[python]{gesibiblio/settings.py, extrait}
EMAIL_BACKEND = "django.core.mail.backends.console.EmailBackend"
DEFAULT_FROM_EMAIL = "GesiBiblio <noreply@gesibiblio.example>"
\end{codebox}

En développement, \code{console.EmailBackend} affiche le courrier dans le terminal.
En production, configurez un vrai serveur SMTP.

\begin{codebox}[html]{templates/registration/password\_reset\_form.html}
{% extends "base.html" %}

{% block titre %}Mot de passe oublié | GesiBiblio{% endblock %}

{% block contenu %}
  <h1>Mot de passe oublié</h1>
  <p>Saisissez votre adresse électronique. Un lien vous sera envoyé.</p>

  <form method="post" class="formulaire">
    {% csrf_token %}
    {{ form.as_p }}
    <button class="bouton bouton-principal" type="submit">
      Envoyer le lien
    </button>
  </form>
{% endblock %}
\end{codebox}

\begin{attention}
La page de demande ne doit pas révéler si l'adresse existe. Django affiche le même
message dans les deux cas. C'est volontaire : cela évite l'énumération des comptes.
\end{attention}

\section{Exercices}

\begin{exercice}[ : page profil protégée]
Vérifiez qu'un visiteur non connecté qui demande \code{/comptes/profil/} est
redirigé vers la page de connexion, puis revient au profil après connexion.
\end{exercice}

\begin{correction}
Ajoutez la route et le décorateur :

\begin{codebox}[python]{catalogue/urls.py, extrait}
path("comptes/profil/", views.profil, name="profil"),
\end{codebox}

\begin{codebox}[python]{catalogue/views.py, extrait}
@login_required
def profil(request):
    ...
\end{codebox}

Avec \code{LOGIN\_URL = "login"}, la redirection est
\code{/comptes/login/?next=/comptes/profil/}. Le gabarit de connexion doit contenir
\code{<input type="hidden" name="next" value="{{ next }}">}.
\end{correction}

\begin{exercice}[ : modification du profil]
Ajoutez une page qui permet de modifier le téléphone et l'adresse du membre
connecté. Utilisez un \code{ModelForm} limité à ces deux champs.
\end{exercice}

\begin{correction}
\begin{codebox}[python]{catalogue/forms.py, extrait}
class ProfilForm(forms.ModelForm):
    class Meta:
        model = Membre
        fields = ["telephone", "adresse"]
\end{codebox}

\begin{codebox}[python]{catalogue/views.py, extrait}
@login_required
def profil_modifier(request):
    membre = request.user.membre

    if request.method == "POST":
        form = ProfilForm(request.POST, instance=membre)
        if form.is_valid():
            form.save()
            messages.success(request, "Votre profil a été mis à jour.")
            return redirect("catalogue:profil")
    else:
        form = ProfilForm(instance=membre)

    return render(
        request,
        "catalogue/profil_modifier.html",
        {"form": form},
    )
\end{codebox}

\begin{note}
Ne demandez jamais le champ \code{user} dans le formulaire. Un utilisateur pourrait
changer le propriétaire du profil. Le formulaire est limité aux champs qu'il a le
droit de modifier.
\end{note}
\end{correction}

\begin{exercice}[ : compteur de connexions]
Affichez sur le profil la date de la dernière connexion et le nombre total de
connexions. Utilisez \code{user.last\_login} et le signal \code{user\_logged\_in}
pour incrémenter un compteur dans le profil.
\end{exercice}

\begin{correction}
\begin{codebox}[python]{catalogue/models.py, extrait}
class Membre(models.Model):
    # ... champs existants
    nombre_connexions = models.PositiveIntegerField(default=0)
\end{codebox}

\begin{codebox}[python]{catalogue/signals.py}
from django.contrib.auth.signals import user_logged_in
from django.dispatch import receiver


@receiver(user_logged_in)
def compter_connexion(sender, request, user, **kwargs):
    if hasattr(user, "membre"):
        user.membre.nombre_connexions += 1
        user.membre.save(update_fields=["nombre_connexions"])
\end{codebox}

Branchez le module dans \code{CatalogueConfig.ready()} :

\begin{codebox}[python]{catalogue/apps.py}
class CatalogueConfig(AppConfig):
    # ... champs existants

    def ready(self):
        from . import signals  # noqa: F401
\end{codebox}

\begin{note}
Le signal s'exécute à chaque connexion. Utilisez \code{update\_fields} pour ne
écrire qu'une colonne. \code{hasattr} évite une erreur pour un superutilisateur
sans profil.
\end{note}
\end{correction}

\begin{retenir}
\begin{itemize}
  \item \code{django.contrib.auth} fournit les utilisateurs, les sessions et les
        permissions.
  \item \code{include("django.contrib.auth.urls")} active les vues standards.
  \item Les gabarits vont dans \code{registration/}.
  \item \code{login\_required} protège une vue fonction ;
        \code{LoginRequiredMixin} protège une classe.
  \item \code{next} ramène l'utilisateur à la page demandée.
  \item La déconnexion se fait en \code{POST} depuis Django 5.
  \item Les mots de passe sont hachés et validés par des validateurs.
\end{itemize}
\end{retenir}
